% !TEX program = pdflatex
% Compilar con pdflatex (2 pasadas) — mismo estilo que Opciones_MHD_2D.tex
\documentclass[11pt,a4paper]{article}

\usepackage[utf8]{inputenc}
\usepackage[T1]{fontenc}
\usepackage[spanish,es-tabla]{babel}
\usepackage[margin=2.3cm]{geometry}
\usepackage{amsmath,amssymb}
\usepackage{booktabs}
\usepackage{enumitem}
\usepackage{xcolor}
\usepackage[colorlinks=true,linkcolor=blue!60!black,citecolor=blue!60!black,urlcolor=blue!60!black]{hyperref}

\newcommand{\B}{\mathbf{B}}
\newcommand{\vct}{\mathbf{v}}
\newcommand{\dif}{\partial_t}
\newcommand{\difdiv}[1]{\nabla\cdot\left(#1\right)}
\newcommand{\code}[1]{\texttt{#1}}

\setlength{\parskip}{2pt}
\setlength{\emergencystretch}{2em}

\title{\textbf{Proyecto 2 — Guía completa: inestabilidad de Rayleigh--Taylor magnetizada}\\[4pt]
\large Todo lo necesario para desarrollar, con perspectiva propia, el problema
\code{Test\_Problems/MHD/Rayleigh\_Taylor} de PLUTO en 2D}
\author{Zeuz Espinosa\\ Astrof\'isica Computacional 2026-II}
\date{Noviembre de 2026}

\begin{document}
\maketitle

\begin{quote}\small
\textbf{Qué es este documento.} Material de referencia unificado para el Proyecto~2:
mapeo de los requisitos del enunciado (\code{project\_2.pdf}), marco teórico completo con
derivaciones (secciones~\ref{sec:fisica}--\ref{sec:nolineal}), especificación exacta del
setup de PLUTO con todos sus parámetros (\S\ref{sec:setup}), teoría lineal con
\emph{predicciones numéricas verificadas} listas para contrastar (\S\ref{sec:lineal}),
plan paso a paso para la simulación PLUTO (\S\ref{sec:pluto}) y para la reproducción
propia en Python (\S\ref{sec:python}), plan del análisis comparativo (\S\ref{sec:comparativo}),
menú de ``perspectivas propias'' rankeadas (\S\ref{sec:perspectivas}) y bibliografía
comentada (\S\ref{sec:biblio}). Los apéndices traen la tabla de números del caso base,
un checklist de entregables y el esqueleto del solver.
\end{quote}

\tableofcontents
\clearpage

%=====================================================================
\section{Requisitos del proyecto y cómo los cubre esta guía}
\label{sec:requisitos}

El enunciado (\code{project\_2.pdf}) pide cinco cosas, cada una con soporte directo aquí:

\begin{table}[htbp]
\centering\small
\caption{Mapeo requisito $\to$ sección de esta guía $\to$ dónde vivirá en el entregable.}
\begin{tabular}{@{}p{5.1cm}p{4.8cm}p{4.9cm}@{}}
\toprule
Requisito (\code{project\_2.pdf}) & Sección de esta guía & Entregable \\
\midrule
1. Problema $\ge$ 2D de \code{Test\_Problems/MHD} & \S\ref{sec:setup}: configuración MHD 2D \#07 (o \#05) del setup \code{Rayleigh\_Taylor} & Archivos PLUTO (\code{definitions.h}, \code{pluto.ini}, \code{init.c}) \\
2. Marco teórico + estado del arte + ecuaciones y fronteras & \S\ref{sec:fisica} (ecuaciones), \S\ref{sec:lineal}--\ref{sec:nolineal} (teoría RT/RT-MHD), \S\ref{sec:contexto} (contexto y estado del arte), \S\ref{sec:biblio} (bibliografía comentada) & Informe: secciones ``Marco teórico'' y ``Estado del arte'' \\
3. Simulación PLUTO completa + gráficas pyPLUTO de todas las variables & \S\ref{sec:pluto} (receta con los comandos de tu flujo local) & Notebook de análisis + figuras \\
4. Reproducción independiente en Python (pipeline completo) & \S\ref{sec:python} (arquitectura del solver, validación escalonada) & Código bien comentado \\
5. Comparación rigurosa: overlays, normas $L_1/L_2$, errores, convergencia & \S\ref{sec:comparativo} (métricas, discrepancias esperadas, plan de convergencia) & Informe: sección ``Análisis comparativo'' \\
\bottomrule
\end{tabular}
\end{table}

Criterios de evaluación del enunciado (marco teórico, corrección PLUTO, calidad Python,
profundidad del análisis, visualización, redacción y citas) se cubren con el checklist del
Apéndice~\ref{app:checklist}. La integridad académica exige citar todo; la bibliografía de
\S\ref{sec:biblio} ya está en formato listo para copiar.

%=====================================================================
\section{Contexto físico y astrofísico}
\label{sec:contexto}

\subsection{La inestabilidad de Rayleigh--Taylor (RT)}

Un fluido \emph{pesado} apoyado sobre uno \emph{ligero} en un campo de gravedad efectiva es
inestable: cualquier corrugación de la interfaz crece, convirtiendo energía potencial
gravitacional en energía cinética. Lord Rayleigh (1883) la estudió para estratificación
continua y G.\ I.\ Taylor (1950) demostró que ocurre siempre que la aceleración apunta del
fluido ligero al pesado (el caso del ``vaso invertido''). El fenómeno es ubicuo: agua sobre
aceite, salmuera en acuíferos, \emph{cavitation}, fusión por confinamiento inercial (ICF),
y una larga lista astrofísica.

El parámetro de densidad que controla todo es el \textbf{número de Atwood}
\begin{equation}
A \;\equiv\; \frac{\rho_\mathrm{pesado}-\rho_\mathrm{ligero}}{\rho_\mathrm{pesado}+\rho_\mathrm{ligero}}
\;\in\; (0,1]\,,
\label{eq:atwood}
\end{equation}
con $A=1$ el límite de burbuja en vacío y $A\to 0$ el límite simétrico. En nuestra
configuración ($\rho_1=1$ abajo, $\rho_2=\eta=2$ arriba): $A=(\eta-1)/(\eta+1)=1/3$.

Mecanismo: al corrugarse la interfaz, el gradiente de presión (vertical, hidrostático) deja
de ser paralelo al gradiente de densidad (normal a la interfaz corrugada); el término
baroclínico $\nabla\rho\times\nabla p / \rho^2$ genera vorticidad del signo que amplifica la
corrugación. El resultado lineal clásico (incompresible, sin viscosidad) es un crecimiento
puramente exponencial con tasa
\begin{equation}
\sigma^2 = A\,g\,k\,,
\qquad k=\frac{2\pi}{\lambda}\,,
\label{eq:rt-hd}
\end{equation}
donde $g$ es la magnitud de la gravedad efectiva y $k$ el número de onda de la perturbación
(§\ref{sec:lineal}). Las longitudes de onda largas crecen más lento ($\sigma\propto\sqrt{k}$);
el corte en onda corta lo pone, en hidrodinámica, la viscosidad o la difusión numérica, y en
nuestro problema, la \emph{tensión magnética}.

En el régimen no lineal la interfaz evoluciona a \emph{burbujas} (fluido ligero subiendo) y
\emph{picos/spikes} (fluido pesado cayendo, con puntas Kelvin--Helmholtz inestables en 3D);
\S\ref{sec:nolineal} resume las leyes de saturación y velocidades.

\subsection{La versión magnetizada: qué cambia y por qué}

Si un campo magnético \emph{horizontal} $\B=B\,\hat{\mathbf{x}}$ permea ambos fluidos
(Kruskal \& Schwarzschild 1954; Chandrasekhar 1961), la línea de campo que cruza la interfaz
corrugada se estira y opone una fuerza restauradora de \textbf{tensión magnética}
$(\B\cdot\nabla)\B/4\pi$ por unidad de volumen. El efecto es anisótropo y selectivo en escala:

\begin{itemize}[nosep]
  \item Modos con vector de onda \emph{paralelo} a $\B$ (``undulares'') sienten la tensión:
  ondas más cortas que la \emph{longitud de corte magnética} $\lambda_c$ quedan
  \textbf{completamente estabilizadas}.
  \item Modos con $\mathbf{k}\perp\B$ (``de intercambio'') no estiran línea alguna y crecen
  como en hidrodinámica pura.
  \item Los modos largos ($\lambda\gg\lambda_c$) crecen casi sin modificaciones, pero el campo
  cambia todo el \emph{régimen no lineal}: suprime las inestabilidades Kelvin--Helmholtz
  secundarias en los flancos de las puntas, alarga y suaviza las burbujas, y se amplifica por
  estiramiento en las puntas de los picos (Stone \& Gardiner 2007).
\end{itemize}

Esto hace del RT magnetizado un problema ideal para un proyecto de MHD computacional: la física
es transparente, el setup es cartesiano con una sola interfaz y un término fuente constante,
hay solución lineal exacta contra la cual medir tasas de crecimiento, y la morfología 2D es
espectacular y muy sensible al tratamiento numérico de $\nabla\cdot\B$ --- exactamente lo que
un análisis comparativo PLUTO vs.\ código propio puede cuantificar.

\subsection{Dónde ocurre en astrofísica (para el estado del arte)}

\begin{itemize}
  \item \textbf{Remanentes de supernova (SNR).} Durante la fase de eyecta dominada, la
  desaceleración del choque externo hace que el contacto acelere \emph{hacia} la eyecta
  densa: RT en la discontinuidad de contacto produce los ``dedos'' observados (Tycho, Cas~A).
  Orlando et al.\ (2012, arXiv:1202.3593) y Fraschetti et al.\ (2010, arXiv:1002.5048)
  muestran en 3D MHD cómo el campo ambiental suprime/modula los dedos y cómo el propio RT
  amplifica campos radiales (Orlando et al.\ 2019). En hidrodinámica 3D: Warren \& Blondin
  (2013).
  \item \textbf{Supernovas Ia.} La ignición y combustión nuclear en una enana blanca es RT
  (llamas de carbono Rayleigh--Taylor); la tasa de crecimiento de las burbujas dicta la
  velocidad de la llama (revisión: Hillebrandt \& Niemeyer 2000).
  \item \textbf{Medio intracúmulo.} Los ``cold fronts'' por oscilamiento (\emph{sloshing})
  de gas frío en cúmulos de galaxias son interfaces RT/KH estabilizadas por campos
  magnéticos drapados; su nitidez observada se usa para acotar $B$ (Markevitch \& Vikhlinin 2007).
  \item \textbf{Prominencias solares.} Los ``plumes'' turbulentos observados por Hinode/SOT
  en quiescentes son RT magnético: en 3D los modos de intercambio crecen alineados con el
  campo y forman flujos ascendentes de sección constante (Hillier et al.\ 2012, 2016;
  Carlyle et al.\ 2017). El régimen no lineal autocoincidente se estudia en Kalluri \&
  Hillier (2025).
  \item \textbf{Laboratorio (HED).} En \emph{magneto-inertial fusion} (MagLIF) y z-pinches,
  la inestabilidad magneto-RT en forros metálicos cilíndricos limita la compresión
  (p.ej.\ Sefkow et al.\ 2014); los experimentos de láser miden las mismas tasas y cortes
  $\lambda_c$ que simulamos aquí.
\end{itemize}

%=====================================================================
\section{Ecuaciones gobernantes}
\label{sec:fisica}

\subsection{MHD ideal con gravedad, en unidades de PLUTO}

PLUTO usa unidades en las que el factor $4\pi$ no aparece (equivalente a racionalizar
$\B\to\B/\sqrt{4\pi}$): la presión total es $p + B^2/2$ y $v_A^2 = B^2/\rho$. Con la gravedad
constante $\mathbf{g}=g\,\hat{\mathbf{y}}$ (en nuestro setup $g=-0.1$) como término fuente, el
sistema es
\begin{align}
\dif\rho + \difdiv{\rho\vct} &= 0\,, \qquad \nabla\cdot\B = 0\,,\\[2pt]
\dif(\rho\vct) + \difdiv{\rho\vct\vct + \Big(p + \tfrac{B^2}{2}\Big)\mathbf{I} - \B\B} &= \rho\,\mathbf{g}\,,\\[2pt]
\dif E + \difdiv{\Big(E + p + \tfrac{B^2}{2}\Big)\vct - (\B\cdot\vct)\B} &= \rho\,\vct\cdot\mathbf{g}\,,\\[2pt]
\dif \B - \nabla\times(\vct\times\B) &= 0\,,
\end{align}
con ecuación de estado ideal
\begin{equation}
E \;=\; \frac{p}{\gamma-1} \;+\; \frac{1}{2}\rho v^2 \;+\; \frac{B^2}{2}\,,
\qquad \gamma = \frac{5}{3}\ \ \text{(valor por defecto, \code{g\_gamma} en \code{Src/globals.h})}.
\end{equation}
(En Gaussianas/CGS habría que reemplazar $B^2/2\to B^2/8\pi$ y $v_A^2=B^2/4\pi\rho$; la
correspondencia es $\B_\mathrm{cgs}=\sqrt{4\pi}\,\B_\mathrm{PLUTO}$ — el \code{init.c} del
problema ya la incorpora, ver \S\ref{sec:setup-Bc}.)

El término fuente en energía $S_E=\rho\,\vct\cdot\mathbf{g}$ es el trabajo gravitacional:
con él, la energía \emph{total} (interna + cinética + magnética + potencial
$\rho\Phi$, con $\mathbf{g}=-\nabla\Phi$) se conserva exactamente en un dominio cerrado ---
éste es un buen diagnóstico global (§\ref{sec:comparativo}).

\subsection{Ondas características y números adimensionales}

Para el CFL y la interpretación física se necesitan las velocidades de onda de la MHD ideal
(magnetosónica rápida, en unidades de PLUTO):
\begin{equation}
\begin{gathered}
c_s^2=\frac{\gamma p}{\rho}\,,
\qquad
v_A^2=\frac{B^2}{\rho}\,,\\[2pt]
c_f^2=\frac{c_s^2+v_A^2}{2}
+\sqrt{\Big(\frac{c_s^2+v_A^2}{2}\Big)^{2}-c_s^2\,v_{A,n}^2}\,,
\end{gathered}
\label{eq:cf}
\end{equation}
con $v_{A,n}$ la componente de la velocidad de Alfvén normal a la dirección del flujo que se
considera. La relación de presión plasmática es $\beta = 2p/B^2$ (CGS: $8\pi p/B^2$).

\subsection{El constraint $\nabla\cdot\B=0$ y sus estrategias numéricas}

La inducción es una ecuación de compatibilidad: si $\nabla\cdot\B=0$ al inicio, la solución
exacta lo conserva. Numéricamente no, y los errores de monopolo producen fuerzas espurias,
transporte erróneo de momento y hasta inestabilidades. Las tres estrategias clásicas
(revisión: Tóth 2000):

\begin{itemize}[nosep]
  \item \textbf{Transporte restringido (CT).} $\B$ vive en caras escalonadas
  ($B_x$ en caras $x$, $B_y$ en caras $y$) y se actualiza con el rotacional de la FEM
  (campo eléctrico inducido $\mathbf{E}=-\vct\times\B$) definido en esquinas;
  $\difdiv{\B}\equiv 0$ \emph{a nivel discretizado}. Es el método de la configuración
  \#07 del setup (y de Athena; Gardiner \& Stone 2005, Londrillo \& Del Zanna 2004).
  Recomendado también para el código propio (\S\ref{sec:python}).
  \item \textbf{Limpieza GLM (Dedner).} Se añade una variable $\psi$ que propaga y amortigua
  los errores de divergencia a velocidad $c_h$: $\dif\B + \nabla\psi = \dots$,
  $\dif\psi + c_h^2\nabla\cdot\B = -\psi/\tau_d$. Más fácil de implementar en malla
  colocalizada; usado por las configuraciones \#06, \#09, \#10 (Mignone et al.\ 2010).
  \item \textbf{Ocho ondas (Powell).} Se añaden términos fuente $\propto\nabla\cdot\B$ que
  restauran conservatividad del sistema extendido; los errores no crecen pero tampoco se
  limpian. PLUTO lo usa solo en variantes particulares.
\end{itemize}

Que el problema tenga configuraciones con CT y con GLM permite una ``meta-comparación''
interesante (§\ref{sec:perspectivas}).

\subsection{Condiciones de frontera del problema RT}

El dominio es $\left[-\tfrac12,\tfrac12\right]\times[-1,1]$ y las fronteras del
\code{pluto.ini} son:

\begin{itemize}[nosep]
  \item \textbf{$x$: periódicas} (la perturbación es un modo único de $\lambda=L_x$).
  \item \textbf{$y$: reflectivas} (paredes rígidas que contienen la atmósfera hidrostática:
  sin flujo de masa ni de energía por arriba ni por abajo).
\end{itemize}

En PLUTO, \code{reflective} es simetría de espejo (verificado en \code{Src/boundary.c}):
en las celdas fantasma se copia reflejado y \emph{cambia de signo la componente normal}
de la velocidad \emph{y} del campo magnético:
\begin{equation}
\text{frontera }y\ \text{(normal }\hat{\mathbf{y}}\text{):}\quad
v_y\to -v_y,\quad B_y\to -B_y,\quad
\rho,p,v_x,B_x\ \text{pares.}
\label{eq:bc}
\end{equation}
Esta receta exacta es la que hay que replicar en Python (\S\ref{sec:python-bc}).

%=====================================================================
\section{El setup \texttt{Rayleigh\_Taylor} de PLUTO, en detalle}
\label{sec:setup}

\subsection{Dominio, equilibrio y normalización}

\begin{table}[htbp]
\centering\small
\caption{Datos geométricos del caso base (configuraciones 2D).}
\begin{tabular}{@{}llp{4.3cm}@{}}
\toprule
Cantidad & Valor & Comentario \\
\midrule
Dominio & $x\in[-0.5,0.5]$, $y\in[-1,1]$ & $L_x=1$ define la unidad de longitud \\
Malla & $256\times512$ (conf.\ 05) o $128\times256$ (conf.\ 07) & uniforme ($\Delta x=\Delta y$) \\
Gravedad $g$ & $-0.1$ & parámetro \code{GRAV}, apunta a $-y$ \\
Densidades & $\rho_1=1$ (abajo, $y\le 0$), $\rho_2=\eta=2$ (arriba) & parámetro \code{ETA}; $A=1/3$ \\
Presión & $p(y)=p_0+\rho\,g\,y$, $\ p_0=1/\gamma=0.6$ & hidrostática exacta \\
$c_s$ (fluido ligero, $y=0^-$) & $1$ & \emph{por diseño}: $p_0$ se elige para esto \\
\bottomrule
\end{tabular}
\end{table}

La condición inicial completa (función \code{Init()} de \code{init.c}) es
\begin{equation}
\begin{gathered}
\rho(y)=\begin{cases}1 & y\le 0\\ \eta & y>0\end{cases},
\qquad
p(y) = \frac{1}{\gamma}+\rho\,g\,y,\\[2pt]
\vct=0\ \ \text{(más la semilla, abajo)},
\qquad
\B = B_x\,\hat{\mathbf{x}}\ \ \text{uniforme}.
\end{gathered}
\label{eq:ci}
\end{equation}
Nótese que $\sigma$ de la teoría lineal no depende de $\gamma$ (la física es incompresible
en la interfaz); $\gamma$ solo fija la velocidad del sonido y el carácter comprimible de la
respuesta.

\subsection{El campo crítico \texorpdfstring{$B_c$}{Bc} y el parámetro \texorpdfstring{$\chi$}{chi}}
\label{sec:setup-Bc}

El \code{init.c} define el \emph{campo magnético crítico} como
\begin{equation}
B_c \;\equiv\; \sqrt{(\rho_2-\rho_1)\,|g|\,L_x}
\;\;\xrightarrow{\ \eta=2,\ L_x=1\ }\;\;
B_c=\sqrt{0.1}\simeq 0.3162\,,
\end{equation}
que en unidades de código (factor $1/\sqrt{4\pi}$) inicializa
\begin{equation}
B_x \;=\; \chi\,\frac{\sqrt{(\eta-1)\,|g|}}{\sqrt{4\pi}}
\;\;\xrightarrow{\ \chi=0.2\ }\;\; B_x \simeq 1.784\times 10^{-2}\,,
\label{eq:bchi}
\end{equation}
donde \code{CHI}$=\chi$ es parámetro de runtime. La definición no es arbitraria: $B_c$ es el
campo que \emph{estabiliza exactamente el modo más largo permitido por el dominio}
($\lambda=L_x$, $k=2\pi/L_x$); la deducción completa está en \S\ref{sec:lineal-bc} y la
verificación numérica en la tabla~\ref{tab:chi}. Con $\chi=0.2$ el campo es dinámicamente
débil en promedio ($v_A\simeq 0.018\ll c_s=1$, $\beta\simeq 3.8\times10^3$) pero decisivo
para las escalas cortas de la interfaz: es toda la gracia del problema.

El campo se inicializa de forma analíticamente divergencia-free mediante el potencial vector
$A_z = y\,B_x$ (\code{ASSIGN\_VECTOR\_POTENTIAL YES}: $\B=\nabla\times\mathbf{A}$ da
$B_x=\partial_y A_z = B_x$ ✓), lo que garantiza un CT limpio desde $t=0$.

\subsection{La perturbación inicial (semilla)}

El setup perturba la velocidad vertical cerca de la interfaz. En 2D, modo único:
\begin{equation}
v_y(x,y) \;=\; -10^{-2}\,\big(1+\cos(2\pi x)\big)\,e^{-200\,y^2}\,.
\label{eq:semilla}
\end{equation}
La envolvente gaussiana tiene ancho $\sim 1/\sqrt{200}\simeq 0.07$; el factor
$(1+\cos 2\pi x)$ es un modo $\lambda=L_x=1$ (más una parte constante en $x$ que solo emite
ondas acústicas débiles y se disipa pronto). La perturbación empuja el fluido pesado
\emph{hacia abajo} en $x=0$ (ahí se formará el pico) y lo levanta en $x=\pm\tfrac12$
(burbujas). Alternativas del setup: perturbación gaussiana 3D, o aleatoria
(\code{USE\_RANDOM\_PERTURBATION YES} en \code{definitions.h}) para sembrar muchos modos.

\subsection{Las 10 configuraciones predefinidas}

\begin{table}[htbp]
\centering\small
\caption{Configuraciones probadas por el autor del setup (cabecera de \code{init.c}).
En negrita las MHD 2D candidatas para este proyecto.}
\begin{tabular}{@{}clllllc@{}}
\toprule
Conf. & Física & Dim & Integrador & Reconstrucción & $\nabla\cdot\B$ & Fuerza \\
\midrule
01 & HD  & 2D & HANCOCK & LINEAR   & --- & vector \\
02 & HD  & 2D & RK3     & MP5\_FD  & --- & vector \\
03 & HD  & 2D & RK3     & PARABOLIC& --- & potencial \\
04 & HD  & 2D & RK2     & LINEAR   & --- & potencial \\
\textbf{05} & \textbf{MHD} & \textbf{2D} & HANCOCK & \textbf{LINEAR}    & \textbf{CT}  & vector \\
\textbf{06} & \textbf{MHD} & \textbf{2D} & RK3     & \textbf{MP5\_FD}   & \textbf{GLM} & vector \\
\textbf{07} & \textbf{MHD} & \textbf{2D} & RK3     & \textbf{PARABOLIC} & \textbf{CT}  & potencial \\
08 & MHD & 2D & ChTr    & PARABOLIC& GLM & potencial \\
09 & MHD & 3D & RK2     & LINEAR   & GLM & vector \\
10 & MHD & 3D & HANCOCK & LINEAR   & GLM & vector \\
\bottomrule
\end{tabular}
\end{table}

\textbf{Recomendación para el proyecto: la \#07 como corrida principal}
(RK3 + PPM + transporte restringido + gravedad por potencial --- fidelidad máxima, la
combinación tipo Athena/Stone \& Gardiner 2007) \textbf{y la \#05 como corrida secundaria}
(más barata: HANCOCK + lineal, ideal para barridos de parámetros). La \#06 (GLM) sirve para
discutir el efecto del control de divergencia. Las tres comparten \code{init.c}, dominio,
fronteras y parámetros \code{ETA=2}, \code{GRAV=-0.1}, \code{CHI=0.2}.

\subsection{\texttt{pluto.ini} de la configuración \#07, anotado}

{\footnotesize
\begin{verbatim}
[Grid]
X1-grid  1  -0.5   128  u   0.5      # x: 128 celdas (Lx = 1)
X2-grid  1  -1.0   256  u   1.0      # y: 256 celdas (Ly = 2, interfaz en y = 0)
X3-grid  1  -0.5    1   u   0.5      # dirección muerta (2D)

[Time]
CFL            0.8      # (0.45 en el ini original; 0.8 corre bien con CT+PPM)
tstop          15.0     # ~6.7 e-foldings nominales; bien entrado el régimen no lineal
first_dt       1.e-3

[Solver]
Solver        roe       # solver de Riemann de Roe

[Boundary]
X1-beg/end    periodic  # modo unico de periodo Lx
X2-beg/end    reflective# paredes rigidas (receta eq. 10)

[Static Grid Output]
dbl           15.0  -1   single_file   # << CAMBIAR a 0.5: ~31 salidas para analisis

[Parameters]
ETA    2.0     # densidad del fluido superior (rho_2)
GRAV  -0.1     # gravedad (debe ser < 0)
CHI    0.2     # campo en unidades del critico Bc
\end{verbatim}
}

La sección \code{[Chombo Refinement]} del ini se ignora en corridas de malla estática
(sin AMR); puede borrarse. Con $\Delta x=1/128$ y $c_f^{\max}\simeq c_s\simeq 1$:
$\Delta t\simeq \mathrm{CFL}\cdot\Delta x \simeq 6\times10^{-3}$, o sea
$\sim 2.5\times10^3$ pasos hasta $t=15$ — minutos en portátil (la \#05 con $256\times512$
toma más, del orden de decenas de minutos).

\subsection{Qué hace cada pieza de \texttt{init.c}}

\begin{itemize}[nosep]
  \item \code{Init()}: llena las primitivas con la eq.~\eqref{eq:ci}, la semilla
  \eqref{eq:semilla} y, si \code{PHYSICS==MHD}, $B_x=\chi B_c/\sqrt{4\pi}$ y
  $A_z=y\,B_x$.
  \item \code{BodyForceVector/Potential}: devuelven $\mathbf{g}=(0,\code{GRAV},0)$ o el
  potencial $\Phi=-\code{GRAV}\,y$. El guard \code{\#ifdef CTU} que invierte $g$ fuera del
  dominio sirve para evaluar consistentemente $\Phi$ en celdas fantasma del esquema
  unsplit CTU.
  \item Parámetros de usuario (\code{USER\_DEF\_PARAMETERS 3}): \code{ETA}, \code{GRAV},
  \code{CHI} (índices 0,1,2), leídos de \code{pluto.ini}.
  \item Constantes de usuario en \code{definitions\_07.h}: \code{ASSIGN\_VECTOR\_POTENTIAL},
  \code{CT\_EMF\_AVERAGE=\allowbreak UCT\_HLL} (promedio de FEM del tipo
  Londrillo--Del Zanna), \code{LIMITER=MC\_LIM}, y
  \code{USE\_RANDOM\_\allowbreak PERTURBATION=\allowbreak NO}.
\end{itemize}

%=====================================================================
\section{Teoría lineal: predicciones cuantitativas verificables}
\label{sec:lineal}

Esta sección es el corazón cuantitativo del proyecto: da números exactos contra los que
medir tanto PLUTO como el código propio.

\subsection{RT hidrodinámica: $\sigma^2=Agk$}

Linealizando las ecuaciones incompresibles alrededor de la interfaz plana con perturbación
$\propto e^{ikx+\sigma t}$, el balance de presión en la interfaz desplazada $\xi$ da la
relación de dispersión clásica $\sigma^2=Agk$ (eq.~\ref{eq:rt-hd}). Con la semilla
\emph{en velocidad} de la eq.~\eqref{eq:semilla}, la evolución lineal exacta es
\begin{equation}
\xi(t) = \xi_0\cosh(\sigma t)+\frac{v_0}{\sigma}\sinh(\sigma t)\,,
\label{eq:xi}
\end{equation}
que con semilla pura en velocidad ($\xi_0=0$) da, para $\sigma t\gtrsim2$,
$\xi(t)\simeq (v_0/2\sigma)\,e^{\sigma t}$,
es decir, una amplitud efectiva inicial $\xi_{0,\mathrm{eff}}=v_0/(2\sigma)$ que luego crece
exponencialmente. Nota: $v_0=0.02$ es el máximo de $|v_y|$ en $x=0$; la gaussiana no es el
autofunction exacto, así que eq.~\eqref{eq:xi} es una aproximación buena (errores de detalle,
no de escala).

\subsection{RT magnetizada: relación de dispersión de Chandrasekhar}
\label{sec:lineal-mhd}

Con $\B=B\,\hat{\mathbf{x}}$ uniforme a ambos lados, la linealización (flujo potencial
incompresible en cada fluido; $\delta\B=(\B\cdot\nabla)\vct/\sigma$ por la inducción) da la
relación de dispersión (Kruskal \& Schwarzschild 1954; Chandrasekhar 1961, §X;
en unidades PLUTO):
\begin{equation}
\boxed{\ \sigma^2 \;=\; A\,g\,k \;-\; \frac{2\,(\mathbf{k}\cdot\B)^2}{\rho_1+\rho_2}\ }
\label{eq:disp}
\end{equation}
(en Gaussianas/CGS, el término magnético se divide entre $4\pi$).
\textbf{Esbozo de deducción} (útil para el informe): para cada fluido el momento linealizado
es $\rho\sigma\vct=-\nabla\delta P + (\B\cdot\nabla)\delta\B$, y la inducción da
$\delta\B=(\B\cdot\nabla)\vct/\sigma$; combinando, la presión total perturbada es
$\delta P=-(\rho\sigma + k^2B^2/\sigma)\phi$ con $\vct=\nabla\phi$, $\phi\propto e^{\mp ky}$.
Imponiendo (i) continuidad de velocidad normal ($\sigma\xi=k\phi_1=-k\phi_2$) y (ii)
continuidad de presión total en la interfaz desplazada (el gradiente hidrostático aporta el
término gravitacional), queda la eq.~\eqref{eq:disp}. El factor 2 es la contribución de
tensión de \emph{ambos} lados de la interfaz. El término magnético es una fuerza restauradora
tipo cuerda: tensión $(\B\cdot\nabla)\delta\B/\sigma^2$ por desplazamiento transversal.

Casos límite que conviene citar en el informe:
\begin{itemize}[nosep]
  \item $\mathbf{k}\perp\B$ (intercambio puro): $\sigma^2=Agk$ --- el campo no hace nada.
  \item $\mathbf{k}\parallel\B$ (undular): estabilización máxima de la eq.~\eqref{eq:disp}.
  En 2D con $\B=B\hat{\mathbf{x}}$ \emph{todo} modo es undular: el setup 2D muestra la
  supresión magnética en su máxima expresión (en 3D los modos de intercambio escapan).
  \item $\B=0$: $\sigma^2=Agk$ ✓.
\end{itemize}

\subsection{Longitud de onda de corte y campo crítico}
\label{sec:lineal-bc}

De la eq.~\eqref{eq:disp}, un modo $\lambda$ con $\mathbf{k}\parallel\B$ es inestable solo si
\begin{equation}
\sigma^2(\lambda) \;=\; A g \frac{2\pi}{\lambda}\Big(1-\chi^2\,\frac{L_x}{\lambda}\Big)\,,
\label{eq:disp-lambda}
\end{equation}
donde se usó $B^2=\chi^2\,(\rho_2-\rho_1)\,g\,L_x/4\pi$ de la eq.~\eqref{eq:bchi}.
de donde sale \emph{todo} lo que sigue (¡verificado numéricamente contra el \code{init.c}!):

\begin{table}[htbp]
\centering\small
\caption{Predicciones lineales del caso base ($\eta=2$, $|g|=0.1$, $L_x=1$, $A=1/3$).
$\sigma_f$ es la tasa del modo fundamental $\lambda=L_x=1$; se cumple exactamente
$\sigma_f=\sigma_\mathrm{HD}\sqrt{1-\chi^2}$ con $\sigma_\mathrm{HD}=0.4576$.}
\begin{tabular}{@{}cccccp{4.35cm}@{}}
\toprule
$\chi$ & $B_x$ (código) & $\lambda_c=\chi^2 L_x$ & $\sigma_f$ & $t_\mathrm{e-fold}=1/\sigma_f$ & Comentario \\
\midrule
0.0 & 0           & ---  & 0.4576 & 2.19 & hidrodinámica pura \\
0.1 & 0.00892     & 0.01 & 0.4553 & 2.20 & \\
0.2 & 0.01784     & 0.04 & 0.4484 & 2.23 & \textbf{caso base del setup} \\
0.3 & 0.02676     & 0.09 & 0.4365 & 2.29 & \\
0.5 & 0.04460     & 0.25 & 0.3963 & 2.52 & \\
0.7 & 0.06244     & 0.49 & 0.3268 & 3.06 & \\
0.9 & 0.08029     & 0.81 & 0.1995 & 5.01 & \\
1.0 & 0.08921     & 1.00 & 0      & ---  & \textbf{marginal: ¡el modo fundamental queda congelado!} \\
\bottomrule
\end{tabular}
\label{tab:chi}
\end{table}

\begin{itemize}[nosep]
  \item \textbf{Longitud de corte:} $\lambda_c=\chi^2 L_x$. Modos con $\lambda<\lambda_c$ a lo
  largo del campo están \emph{estabilizados}. Con $\chi=0.2$: $\lambda_c=0.04$.
  \item \textbf{Campo crítico:} $\chi=1$ ($B=B_c$) estabiliza \emph{exactamente} el modo más
  largo del dominio ($\lambda=L_x$); $\chi>1$ estabiliza todo (por eso el \code{init.c}
  llama ``crítico'' a $B_c$; Boyd \& Sanderson 2003, p.~99).
  \item \textbf{Modo de máximo crecimiento:} derivando \eqref{eq:disp-lambda},
  $\lambda_\mathrm{max}=2\lambda_c$ con $\sigma_\mathrm{max}
  =\sqrt{A g \pi/(2\chi^2 L_x)}$. Con $\chi=0.2$: $\lambda_\mathrm{max}=0.08$,
  $\sigma_\mathrm{max}\simeq1.14$ --- pero \emph{ojo}: ese modo no está sembrado por la CI;
  solo aparece por inestabilidades secundarias/ruido numérico. Físicamente: el corte
  magnético limita el crecimiento a la tasa de cruce de Alfvén de la escala,
  $\sigma_\mathrm{max}\sim v_A/\lambda_c$ (compruébese).
  \item \textbf{El fundamental del setup} ($\lambda=1\gg\lambda_c=0.04$) apenas nota el
  campo: $\sigma$ cae $2\%$. Toda la acción magnética está en las escalas cortas que
  aparecen al volverse no lineal el interfaz (§\ref{sec:nolineal}).
\end{itemize}

\subsection{Límites de validez de la teoría}

La eq.~\eqref{eq:disp} supone fluidos incompresibles, invíscidos, semi-infinitos y campo
uniforme. En nuestro setup: (i) el número de Mach de la fase lineal es
$\sigma/(k c_s)\simeq0.07$ — correcciones compresibles $\lesssim$ pocos por ciento;
(ii) la viscosidad física es nula, pero la \emph{numérica} del esquema corta los modos
con $\lambda\sim$ pocas celdas: en la malla \#07 ($\Delta x=1/128$), $\lambda_c=0.04$
equivale a $\sim5$ celdas --- marginalmente resuelto; en la \#05 ($\Delta x=1/256$), $\sim10$
celdas. \emph{Consecuencia práctica}: la fidelidad del corte magnético exige resolución;
discutir esto en el informe es análisis de convergencia gratis (§\ref{sec:comparativo}).

\subsection{Cronología esperada de la corrida base ($\chi=0.2$)}

Con $\sigma=0.4484$ y $\xi_{0,\mathrm{eff}}=v_0/(2\sigma)\simeq0.022$ (eq.~\ref{eq:xi},
$v_0=0.02$):
\begin{table}[htbp]
\centering\small
\caption{Hitos estimados del caso base (la no linealidad real llega un poco antes/así
por la geometría del hongo).}
\begin{tabular}{@{}cl@{}}
\toprule
$t$ & Evento (estimación) \\
\midrule
$\sim$1.8 & $\xi\simeq0.05$: interfaz claramente deformada \\
$\sim$3.4 & $\xi\simeq0.1\,\lambda$: inicio del régimen no lineal \\
$\sim$6.9 & $\xi\simeq\lambda/2$: burbujas/picos desarrollados \\
$\gtrsim$8 & saturación, picos KH/plasmoides, mezcla \\
15 (tstop) & régimen no lineal tardío \\
\bottomrule
\end{tabular}
\end{table}

%=====================================================================
\section{Régimen no lineal: burbujas, picos y el papel del campo}
\label{sec:nolineal}

\subsection{Sin campo: Layzer y compañía}

Cuando $\xi\sim\lambda$, la fase exponencial muere y la burbuja asciende a velocidad
\emph{constante}. El modelo de Layzer (1955), generalizado a $A$ arbitrario por Goncharov
(2002), da en 2D
\begin{equation}
v_b \;=\; \sqrt{\frac{2A\,g\,\lambda}{\pi\,(1+A)}}
\;\;\xrightarrow{\ A=1/3,\ g=0.1,\ \lambda=1\ }\;\; v_b \simeq 0.126\,,
\label{eq:layzer}
\end{equation}
en unidades de código,
y el pico cae más rápido cuanto mayor $A$ (límite $A\to1$: caída libre
$v_s\to\sqrt{2\,g\,\xi}$). Para $A=1/3$ ambos son comparables, $v_s\gtrsim v_b$.
Estas velocidades asintóticas son medibles en las simulaciones (velocidad del tope de la
burbuja / punta del pico) y dan otra verificación independiente de ambos códigos.

\subsection{Con campo horizontal (2D): lo que encontró Stone \& Gardiner (2007)}

La referencia central del proyecto. Para 2D con $\mathbf{k}\parallel\B$ (nuestro caso) y
campos sub-críticos ($\chi<1$):

\begin{enumerate}[nosep]
  \item \textbf{El crecimiento del modo fundamental es esencialmente el hidrodinámico}
  (nuestra tabla~\ref{tab:chi}: $-2\%$ en $\sigma$ para $\chi=0.2$).
  \item \textbf{Supresión de las inestabilidades secundarias.} Los flancos de las puntas
  son interfaces Kelvin--Helmholtz cuya escala más inestable es corta ($\lambda\lesssim$
  espesor de corte); el campo estabiliza todo $\lambda<\lambda_c$: los hongos pierden los
  ``rulos'' KH, las puntas se vuelven anchas y alargadas, las burbujas se ensanchan.
  \item \textbf{Amplificación del campo por estiramiento:} el flujo estira las líneas en las
  puntas de los picos; $|\B|_\mathrm{max}/B_0$ crece con el tiempo y la tensión se vuelve
  dinámicamente importante tarde en la evolución (el ``campo se desquita'' de ser débil).
  \item \textbf{Plasmoides/islas magnéticas} se forman a lo largo de la interfaz
  (corrientes sheet fragmentadas) cuando la corriente se concentra en las puntas.
  \item \textbf{Alineación:} la tensión tiende a alinear la interfaz con el campo, frenando
  la mezcla tardía comparada con el caso HD.
\end{enumerate}
Todo esto es cualitativa y cuantitativamente chequeable en las figuras de PLUTO/Python
(morfología, $|\B|$ máximo en el tiempo, tamaño mínimo de estructura
$\sim\lambda_c$) y es el contenido físico del análisis comparativo.

\subsection{2D vs.\ 3D}

En 3D con el mismo $\B$ horizontal, los modos con $\mathbf{k}\perp\B$ crecen sin freno
magnético: la mezcla es mucho más vigorosa y turbulenta que en 2D (donde \emph{todo} modo es
undular y está frenado). El setup 2D del proyecto no es una ``fotografía'' del caso 3D sino
el caso de supresión máxima; decir esto con propiedad (y citar a Stone \& Gardiner 2007 y
Kalluri \& Hillier 2025 para el régimen no lineal 3D) suma puntos de estado del arte.

\subsection{Energética (diagnóstico global)}

En el dominio cerrado (fronteras reflectivas) con el término fuente gravitacional completo:
\begin{equation}
\underbrace{\frac{d}{dt}\int\!\Big(\tfrac{p}{\gamma-1}+\tfrac{\rho v^2}{2}+\tfrac{B^2}{2}\Big)dV}_{E_\mathrm{tot\,int+cin+mag}}
=\;-\,\frac{d}{dt}\int \rho\,\Phi\,dV
\;\;\Longleftrightarrow\;\;
\Delta E_\mathrm{cin}+\Delta E_\mathrm{mag}= -\Delta E_\mathrm{pot}-\Delta E_\mathrm{int}\,.
\end{equation}
La inestabilidad convierte $E_\mathrm{pot}$ en cinética (y algo en magnética por
estiramiento); registrar $E_\mathrm{cin}(t)$, $E_\mathrm{mag}(t)$, $E_\mathrm{pot}(t)$ en
ambos códigos es un diagnóstico global robusto, menos sensible al caos de interfaz que las
comparaciones puntuales (§\ref{sec:comparativo}).

%=====================================================================
\section{Simulación con PLUTO: receta paso a paso}
\label{sec:pluto}

Flujo validado en tu equipo (Fedora, gcc, PLUTO 4.4 serial, pyPLUTO 1.2.6; ver tu
\code{COMANDOS\_\allowbreak{}\_\allowbreak PLUTO\_\allowbreak{}\_\allowbreak CP\_\allowbreak{}\_\allowbreak Alfven.md}).
Con \code{\$PLUTO\_DIR} apuntando a
\code{.../Astrof\'isica\_Computacional/PLUTO} (¡ruta con guion bajo!):

\subsection{Preparar la corrida}

{\footnotesize
\begin{verbatim}
cd ~/Documents/UN/2026-2/Astrofísica_Computacional/Proyecto_2
mkdir rt_mhd && cd rt_mhd
cp "$PLUTO_DIR"/Test_Problems/MHD/Rayleigh_Taylor/{init.c,definitions_07.h,pluto_07.ini} .
cp definitions_07.h definitions.h      # setup.py lee definitions.h
cp pluto_07.ini       pluto.ini
# salida para el analisis: 31 cuadros en vez de solo el final
sed -i 's|^dbl .*|dbl        0.5  -1   single_file|' pluto.ini
\end{verbatim}
}

\subsection{Compilar y ejecutar}

{\footnotesize
\begin{verbatim}
printf 'ARCH = Linux.gcc.defs\n' > makefile
python3 "$PLUTO_DIR/setup.py" --auto-update   # genera makefile completo y enlaza Src/
echo 'CFLAGS += -D_DEFAULT_SOURCE' >> makefile # fix srand48 en Fedora
make && ./pluto          # minutos: 128x256, RK3, t=15
\end{verbatim}
}

Reglas de oro (ya aprendidas en la tarea de Alfvén): tras editar \code{definitions.h}
$\to$ \code{make clean} y regenerar el makefile (\code{setup.py --auto-update} +
re-añadir \code{CFLAGS}); cambios de \code{pluto.ini} no requieren recompilar. Útiles:
\code{./pluto -maxsteps N}, \code{./pluto -restart k}.

Para el barrido en $\chi$ (§\ref{sec:perspectivas}): duplicar la carpeta por valor de
$\chi$ y ajustar \code{CHI} en \code{pluto.ini} (o usar \code{./pluto -restart 0 -pars
"CHI:0.5"} si tu versión soporta \code{-pars}).

\subsection{Postproceso con pyPLUTO 1.2.6}

{\footnotesize
\begin{verbatim}
import pyPLUTO as pp, numpy as np, matplotlib.pyplot as plt
D = pp.Load(30, path='.')            # salida t=15; D.ntime da el tiempo
rho, p   = D.rho, D.prs
vx, vy   = D.vx1, D.vx2
Bx, By   = D.Bx1, D.Bx2
x, y     = D.x1, D.x2
# mapa de densidad estilo informe:
plt.figure(figsize=(5,8))
plt.pcolormesh(x, y, rho.T, shading='auto', cmap='viridis')
plt.colorbar(label=r'$\rho$'); plt.xlabel('x'); plt.ylabel('y')
plt.title(f't = {D.ntime:.1f}'); plt.tight_layout(); plt.savefig('rho_t15.png', dpi=150)
\end{verbatim}
}

\textbf{Figuras obligatorias por el enunciado} (todas las variables físicas, con
colorbar/ejes/título): $\rho$, $p$, $v_x$, $v_y$, $|\vct|$, $B_x$, $B_y$, $|\B|$
— mismas convenciones para PLUTO y Python para que el overlay sea directo. Sugeridas
adicionales: vorticidad $\omega_z=\partial_x v_y-\partial_y v_x$ (los ``rulos'' KH),
$|\B|/B_0$ (amplificación), y la posición de la interfaz $y_\mathrm{int}(x,t)$ (contorno
$\rho=1.5$).

%=====================================================================
\section{Reproducción independiente en Python}
\label{sec:python}

\subsection{Arquitectura recomendada}

Para que la comparación con PLUTO \#07 (CT) sea de tú a tú, el diseño natural es
\textbf{volúmenes finitos + CT escalonado + método de líneas}:

\begin{itemize}[nosep]
  \item \textbf{Estado}: conservativas en celdas
  $U=(\rho,\ \rho u,\ \rho v,\ E)$ \emph{colocadas}, y el campo en caras:
  $B_x$ en caras $x$ ($i\pm\tfrac12,j$), $B_y$ en caras $y$ ($i,j\pm\tfrac12$).
  \item \textbf{Reconstrucción}: lineal con limitador MC en variables primitivas
  (como PLUTO \#05; PPM si sobra tiempo).
  \item \textbf{Flujos}: solver de Riemann \textbf{HLL} (robusto y corto de programar;
  HLLD si se quiere fidelidad extra). Para las caras se necesita $B$ reconstruido en la
  cara (promedio de las dos celdas adyacentes a cada lado, o reconstrucción directa).
  \item \textbf{Inducción por CT}: FEM $E_z=(\vct\times\B)_z=v_x B_y-v_y B_x$ en esquinas
  $(i\pm\tfrac12,j\pm\tfrac12)$, promediando las FEM de las 4 caras adyacentes
  (``corner average''; Gardiner \& Stone 2005; PLUTO usa la variante más precisa
  UCT\_HLL de Londrillo \& Del Zanna 2004, no hace falta igualarla).
  \item \textbf{Tiempo}: RK2 (o RK3) no dividido; el paso de actualización de $\B$ va
  dentro de cada etapa RK con la $E_z$ de esa etapa.
  \item \textbf{Fuente}: $\rho g$ en momento y $\rho\vct\cdot\mathbf{g}$ en energía.
\end{itemize}

Actualización CT en 2D (mismas fórmulas para el código y para entender PLUTO):
\begin{equation}
\frac{\partial B_x}{\partial t} = -\frac{\partial E_z}{\partial y}\Big|_\text{cara x},
\qquad
\frac{\partial B_y}{\partial t} = +\frac{\partial E_z}{\partial x}\Big|_\text{cara y},
\qquad E_z = v_x B_y - v_y B_x\,,
\label{eq:ct}
\end{equation}
discretizado con diferencias centradas entre caras/esquinas vecinas: $\nabla\cdot\B$
numérico queda \emph{idénticamente} cero a máquina.

Flujo HLL en la dirección $x$ (ídem $y$):
\begin{equation}
\mathbf{F}_{i+1/2} =
\frac{S_R\,\mathbf{F}_L - S_L\,\mathbf{F}_R + S_L S_R\,(\mathbf{U}_R-\mathbf{U}_L)}{S_R-S_L}\,,
\qquad
\begin{cases}
S_L=\min(u_L-c_{f,L},\,u_R-c_{f,R})\\
S_R=\max(u_L+c_{f,L},\,u_R+c_{f,R})
\end{cases}
\label{eq:hll}
\end{equation}
con $c_f$ de la eq.~\eqref{eq:cf}. El flujo HLL de la inducción ($-E_z$) \emph{no} se usa para
actualizar $\B$ (eso lo hace CT); sí sirve como $E_z$ de cara para el promedio de esquinas.

\subsection{Paso de tiempo (CFL)}

\begin{equation}
\Delta t = C_\mathrm{CFL}\,
\min\!\Big(\frac{\Delta x}{\max(|u|+c_{f,x})},\ \frac{\Delta y}{\max(|v|+c_{f,y})}\Big),
\qquad C_\mathrm{CFL}\simeq 0.3\text{--}0.8\ \ (\text{RK2+CT tolera } \lesssim 0.5).
\end{equation}

\subsection{Condiciones de frontera en el código propio}
\label{sec:python-bc}

Periódicas en $x$ (roll de arrays). Reflectivas en $y$ según la eq.~\eqref{eq:bc}:
en las celdas fantasma, copia espejada con $v_y\to-v_y$ y $B_y\to-B_y$
($B_y$ vive en caras $y$: la cara exactamente en la pared vale $0$ y las caras fantasma se
copian con signo invertido respecto a su imagen). $\rho$, $p$, $v_x$, $B_x$, $E_z$
(evaluada con el $v_y$ invertido) pares.

\subsection{El equilibrio hidrostático: \emph{well-balancing} casero}

El peor error sistemático en RT numérico es romper el balance $-\nabla p+\rho\mathbf{g}=0$
en la interfaz: aparecen velocidades espurias $\sim$ error de truncamiento que contaminan la
fase lineal. Mitigaciones (en orden de costo):

\begin{enumerate}[nosep]
  \item \textbf{Inicialización discreta consistente:} en vez de $p_j=p_0+\rho g y_j$,
  integrar celda a celda $p_{j}=p_{j-1} - \tfrac12(\rho_j+\rho_{j-1})\,|g|\,\Delta y$
  (reconstruye la interfaz $\rho$ promediada en la celda que la corta).
  \item \textbf{Test de equilibrio:} correr \emph{sin perturbación}; exigir
  $\max|v|<10^{-3}$ (idealmente $\lesssim10^{-4}$) hasta $t=10$ antes de confiar en el solver.
  \item (Opcional, avanzado) evolucionar perturbaciones $q'=q-q_\mathrm{hidro}(y)$ en un
  esquema \emph{well-balanced} exacto.
\end{enumerate}

\subsection{Escalera de validación (¡en orden!)}

\begin{table}[htbp]
\centering\small
\caption{Tests incrementales con criterio de éxito. Cada paso reusa el anterior.}
\begin{tabular}{@{}clp{7.2cm}@{}}
\toprule
\# & Test & Criterio de éxito \\
\midrule
1 & Equilibrio hidrostático sin semilla & $\max|v|<10^{-3}$ hasta $t=10$ \\
2 & Tubo de choque MHD 1D (Ryu \& Jones 1a) & perfiles dentro de ruido de gráfico vs.\ referencia \\
3 & Onda de Alfvén 2D diagonal (¡ya la tienes en 1D!) & $\sigma$ numérico $<$1\% de error; $\nabla\cdot\B\approx10^{-15}$ \\
4 & RT hidrodinámico ($\chi=0$) & $\sigma$ medido dentro de $\sim$5\% de $0.4576$ (ajuste $3\lesssim t\lesssim6$) \\
5 & RT MHD $\chi=0.2$ & $\sigma$ dentro de $\sim$5--8\% de $0.4484$ \\
6 & Convergencia & $L_1(\rho)$ vs.\ $\Delta x$ en $t=3$ con pendiente $\sim$1--2 (§\ref{sec:comparativo}) \\
\bottomrule
\end{tabular}
\end{table}

\subsection{Alternativa GLM (si el CT escalonado se vuelve cuesta arriba)}

Sistema extendido colocado (todo en celdas), más fácil de programar:
$\partial_t\B+\nabla\psi=-\nabla\times(\vct\times\B)$,
$\partial_t\psi+c_h^2\nabla\cdot\B=-\psi/\tau_d$,
con $c_h=\max$ velocidad característica del paso y $\tau_d\sim$ pocos $\Delta t$
(Mignone, Tzeferacos \& Bodo 2010). Es la opción de las configs \#06/\#09/\#10.
Costo: $\nabla\cdot\B$ residual $\neq0$ (típ.\ $10^{-4}$--$10^{-3}$ relativo) que se debe
reportar como parte del análisis comparativo.

\subsection{Resolución y costo}

Recomendado: \textbf{Python a $128\times256$} (igual que PLUTO \#07) con corridas extra a
$64\times128$ y $256\times512$ para convergencia. Con numpy vectorizado, el caso base son
$\sim$2--3~mil pasos: minutos a decenas de minutos por corrida. Numba/JAX si hay barridos.
El esqueleto de clases está en el Apéndice~\ref{app:esqueleto}.

%=====================================================================
\section{Análisis comparativo PLUTO vs.\ Python: plan}
\label{sec:comparativo}

\subsection{Métricas}

\begin{itemize}[nosep]
  \item \textbf{Normas} (en cuadrícula común, interpolando la más fina a la más gruesa;
  normalizadas para adimensionalizar):
  \begin{equation}
  \begin{gathered}
  L_1(q) = \frac{1}{A_\mathrm{dom}}\int |q_\mathrm{PLUTO}-q_\mathrm{py}|\,dA,
  \qquad
  L_2(q) = \Big(\frac{1}{A_\mathrm{dom}}\int (q_\mathrm{PLUTO}-q_\mathrm{py})^2\,dA\Big)^{1/2},\\[2pt]
  \epsilon_\mathrm{rel} = \frac{L_1(q)}{\langle|q_\mathrm{PLUTO}|\rangle}\,.
  \end{gathered}
  \end{equation}
  \item \textbf{Interfaz:} $y_\mathrm{int}(x,t)$ = contorno $\rho=1.5$; amplitud
  $a(t)=\tfrac12[\max y_\mathrm{int}-\min y_\mathrm{int}]$; ajuste lineal de $\ln a$ vs $t$
  en $3\lesssim t\lesssim6$ $\to$ $\sigma_\mathrm{PLUTO}$, $\sigma_\mathrm{py}$ vs
  $\sigma_\mathrm{teoría}=0.4484$.
  \item \textbf{Energías:} $E_\mathrm{cin}(t)$, $E_\mathrm{mag}(t)$, $E_\mathrm{pot}(t)$
  (integrales discretas idénticas en ambos códigos) --- comparación robusta al caos.
  \item \textbf{Cortes 1D:} $\rho(y)$ en $x=0$ (eje del pico) y $x=\pm0.25$ (frente de
  burbuja) en $t=\{5,10,15\}$, superpuestos.
  \item \textbf{Divergencia:} $\max|\nabla\cdot\B|\,\Delta x/B_0$ vs $t$ para ambos
  (CT: $\sim10^{-15}$; GLM: $\sim10^{-4}$).
\end{itemize}

\subsection{Gráficas del informe (lista mínima)}

Mapas 2D de $\rho$, $|\vct|$, $|\B|$ (PLUTO | Python lado a lado y/o overlay por
diferencia $q_\mathrm{PLUTO}-q_\mathrm{py}$) en $t=\{3,6,10,15\}$; $a(t)$ en escala
semilog con las tres pendientes (PLUTO, Python, teoría); cortes 1D; energías;
$\sigma$ vs $\chi$ (si se hace el barrido, §\ref{sec:perspectivas}); tabla de normas
$L_1/L_2$ por variable y tiempo.

\subsection{Discrepancias esperadas (para discutir con anticipación)}

\begin{table}[htbp]
\centering\small
\caption{Fuentes de discrepancia PLUTO\#07 vs.\ código propio, y su firma.}
\begin{tabular}{@{}p{4.6cm}p{4.4cm}p{6.6cm}@{}}
\toprule
Fuente & Firma en los datos & Mitigación/discusión \\
\midrule
Solver (Roe vs.\ HLL) y reconstrucción (PPM vs.\ lineal+MC) & Interfaces PLUTO más nítidas; $\sigma_\mathrm{py}$ levemente menor; más difusión en $|\nabla\rho|$ & Es la discusión central: disipación numérica $\equiv$ viscosidad efectiva que corta ondas cortas \\
Resolución & Estructuras sub-$\lambda_c$ solo a alta resolución & Correr 64/128/256 y reportar \\
Frontera/salida en tiempos distintos & Fases corridas en overlays & Fijar \code{dbl 0.5} y \code{tstop} idénticos; interpolar a $t$ comunes \\
Tratamiento de $\nabla\cdot\B$ (CT vs.\ GLM) & Diferencias en tensión cerca de puntas/plasmoides & Reportar $\nabla\cdot\B$ de cada código \\
Balance hidrostático imperfecto (Python) & $\sigma$ contaminado en $t\lesssim2$ & Test de equilibrio previo (\S\ref{sec:python}) \\
Sensibilidad caótica tardía & A partir de $t\sim$7--8, las posiciones puntuales divergen entre \emph{cualquier} par de corridas & Comparar estadísticas (energías, espectros) en vez de punto a punto \\
\bottomrule
\end{tabular}
\end{table}

\textbf{Punto clave para el informe}: la comparación punto a punto es limpia solo en la fase
lineal y temprana no lineal ($t\lesssim6$); después la RT es caótica (dos corridas de PLUTO
con semillas trivialemente distintas divergen) y lo honesto es comparar estadística y
morfología. Decirlo explícitamente demuestra entendimiento.

\subsection{Convergencia}

Correr el código propio (y PLUTO si se quiere: \#05 a $256\times512$ vs \#07 a $128\times256$
vs re-mesh de \#07) en $N_x=\{64,128,256\}$. Medir $L_1(\rho)$ entre resoluciones
consecutivas a $t=3$ (aún determinista): en escala log-log la pendiente es el orden efectivo
$\sim$1 (lineal+MC) a $\sim$2 (PPM). Reportar también cómo $\sigma$ medido converge a
$0.4484$ con $N_x$.

%=====================================================================
\section{Perspectivas propias: menú rankeado}
\label{sec:perspectivas}

Para ``darle perspectiva propia'' (y diferenciarse), en orden de relación impacto/esfuerzo:

\begin{enumerate}
  \item \textbf{Barrido en $\chi$ contra la teoría.}
  $\chi\in\{0,\,0.1,\,0.2,\,0.3,\,0.5,\,0.7,\,0.9,\,1.0\}$ en PLUTO \#05 (barata) y en
  Python; medir $\sigma(\chi)$ y superponer la predicción
  $\sigma=\sigma_\mathrm{HD}\sqrt{1-\chi^2}$ (tabla~\ref{tab:chi}). El punto $\chi=1$
  (interfaz congelada) y el corte morfológico $\lambda_c=\chi^2$ son resultados
  espectaculares y cuantitativos. \emph{Este es el experimento insignia.}
  \item \textbf{Atwood:} $\eta\in\{1.5,2,3,10\}$; verificar $\sigma_\mathrm{HD}\propto\sqrt{A}$
  y la velocidad de burbuja de Layzer $v_b\propto\sqrt{A/(1+A)}$ (eq.~\ref{eq:layzer})
  midiendo el tope de la burbuja en $t$ grande.
  \item \textbf{Energía y amplificación:} $|\B|_\mathrm{max}(t)/B_0$ (estiramiento en las
  puntas, §\ref{sec:nolineal}) y partición $E_\mathrm{cin}$/$E_\mathrm{mag}$ vs $t$
  para varios $\chi$; discutir cuándo el campo se vuelve dinámicamente importante
  (Stone \& Gardiner 2007).
  \item \textbf{Semilla aleatoria} (\code{USE\_RANDOM\_PERTURBATION YES}): mezcla multimodo;
  medir el tamaño mínimo de estructura vs $\lambda_c$ (el ``cutoff magnético'' visible).
  \item \textbf{Meta-numérica:} PLUTO \#05 vs \#06 vs \#07 (CT vs GLM, lineal vs PPM):
  cómo cambia $\nabla\cdot\B$, la nitidez de la interfaz y $\sigma$. Es análisis comparativo
  ``interno'' que enriquece la discusión de métodos.
  \item \textbf{(Ambicioso) 3D:} conf.\ \#09/\#10 de PLUTO vs 2D: los modos de intercambio
  $\perp\B$ escapan a la supresión (§\ref{sec:nolineal}); con solo PLUTO (sin reproducir 3D
  en Python) ya es un bonito apéndice.
\end{enumerate}

%=====================================================================
\section{Bibliografía comentada}
\label{sec:biblio}

\paragraph{Teoría clásica.}
Rayleigh (1883) y Taylor (1950) son el origen HD; Kruskal \& Schwarzschild (1954) añaden el
campo; \textbf{Chandrasekhar (1961), caps.~X} es LA referencia de la relación de dispersión
(eq.~\ref{eq:disp}) — citarla para la teoría lineal. Sharp (1984) y Kull (1991) son las
revisiones clásicas del RT HD y de sus aplicaciones; Layzer (1955), Davies \& Taylor (1950) y
Goncharov (2002) cubren el régimen no lineal de burbujas (eq.~\ref{eq:layzer}).

\paragraph{Métodos numéricos / MHD.}
\textbf{Stone \& Gardiner (2007)} es el artículo central del proyecto: RT MHD 2D/3D con CT,
supresión de KH secundaria, amplificación de campo — con quien comparar morfología y
$\sigma$. Gardiner \& Stone (2005) y Londrillo \& Del Zanna (2004) documentan el CT unsplit
que usa PLUTO \#07; Tóth (2000) revisa todas las estrategias $\nabla\cdot\B$; Mignone et al.\
(2010) el GLM de \#06; Mignone et al.\ (2007, 2012) son los papers del código PLUTO.
Dimonte et al.\ (2004) —la comparación ``$\alpha$-group'' de códigos RT— es el precedente
metodológico perfecto para justificar tu análisis comparativo.

\paragraph{Astrofísica y laboratorio (estado del arte).}
SNR magnetizados: Orlando et al.\ (2012; y 2019 para campos radiales), Fraschetti et al.\
(2010), Warren \& Blondin (2013, HD 3D). SNe Ia: Hillebrandt \& Niemeyer (2000).
Intracúmulo: Markevitch \& Vikhlinin (2007). Prominencias solares: Hillier et al.\
(2012, 2016), Carlyle et al.\ (2017) y Kalluri \& Hillier (2025, régimen no lineal
autocoincidente). Laboratorio HED: Sefkow et al.\ (2014). El ``campo crítico'' del
\code{init.c} sigue a Boyd \& Sanderson (2003, p.~99).

\begin{thebibliography}{99}\small

\bibitem{rayleigh1883} Rayleigh, Lord 1883, ``Investigation of the character of the equilibrium of an incompressible heavy fluid of variable density'', Proc.\ London Math.\ Soc.\ 14, 170.
\bibitem{taylor1950} Taylor, G.~I. 1950, ``The instability of liquid surfaces when accelerated in a direction perpendicular to their planes'', Proc.\ R.\ Soc.\ Lond.\ A 201, 192.
\bibitem{kruskal1954} Kruskal, M. \& Schwarzschild, M. 1954, ``Some instabilities of a completely ionized plasma'', Proc.\ R.\ Soc.\ Lond.\ A 223, 348.
\bibitem{chandra1961} Chandrasekhar, S. 1961, \emph{Hydrodynamic and Hydromagnetic Stability} (Dover), caps.~X (RT magnetizado, §§89--101).
\bibitem{daviestaylor1950} Davies, R.~M. \& Taylor, G.~I. 1950, ``The mechanics of large bubbles rising through extended liquids...'', Proc.\ R.\ Soc.\ Lond.\ A 200, 375.
\bibitem{layzer1955} Layzer, D. 1955, ``The stability of accelerated fluid sheets'', ApJ 122, 1.
\bibitem{sharp1984} Sharp, D.~H. 1984, ``An overview of Rayleigh--Taylor instability'', Physica D 12, 3.
\bibitem{kull1991} Kull, H.-J. 1991, ``Theory of the Rayleigh--Taylor instability'', Phys.\ Rep.\ 206, 197.
\bibitem{goncharov2002} Goncharov, V.~N. 2002, ``Analytical model of nonlinear, single-mode, classical Rayleigh--Taylor instability at arbitrary Atwood numbers'', Phys.\ Rev.\ Lett.\ 88, 134502.
\bibitem{jun1995} Jun, B.-I., Norman, M.~L. \& Stone, J.~M. 1995, ``A numerical study of Rayleigh--Taylor instability in magnetic fluids'', ApJ 453, 332.
\bibitem{sg2007} Stone, J.~M. \& Gardiner, T.~A. 2007, ``Nonlinear evolution of the magnetohydrodynamic Rayleigh--Taylor instability'', Phys.\ Fluids 19, 094104 (arXiv:0707.1022).
\bibitem{gs2005} Gardiner, T.~A. \& Stone, J.~M. 2005, ``An unsplit Godunov method for ideal MHD via constrained transport'', JCP 205, 509.
\bibitem{ld2004} Londrillo, P. \& Del Zanna, L. 2004, ``High-order upwind schemes for multidimensional MHD'', JCP 195, 17.
\bibitem{toth2000} T\'oth, G. 2000, ``The $\nabla\cdot\B=0$ constraint in shock-capturing MHD codes'', JCP 161, 605.
\bibitem{mtb2010} Mignone, A., Tzeferacos, P. \& Bodo, G. 2010, ``High-order conservative schemes for MHD... GLM'', JCP 229, 2117.
\bibitem{pluto2007} Mignone, A., Bodo, G., Massaglia, S. et al. 2007, ``PLUTO: a numerical code for computational astrophysics'', ApJS 170, 228.
\bibitem{pluto2012} Mignone, A., Zanni, C., Tzeferacos, P. et al. 2012, ``The PLUTO code for adaptive mesh computations...'', ApJS 198, 7.
\bibitem{dimonte2004} Dimonte, G. et al. 2004, ``A comparative study of the numerical simulation of the Rayleigh--Taylor instability'', Phys.\ Fluids 16, 1668.
\bibitem{orlando2012} Orlando, S. et al. 2012, ``3D MHD simulations of magnetized supernova remnant'' (arXiv:1202.3593).
\bibitem{fraschetti2010} Fraschetti, F. et al. 2010, ``Simulation of the growth of the 3D Rayleigh--Taylor instability in young supernova remnants'' (arXiv:1002.5048).
\bibitem{warren2013} Warren, D.~C. \& Blondin, J.~M. 2013, ``The dependence of Rayleigh--Taylor instability on the structure of a Type Ia supernova'', MNRAS 429, 3099.
\bibitem{hillebrandt2000} Hillebrandt, W. \& Niemeyer, J.~C. 2000, ``Type Ia supernova explosion models'', ARA\&A 38, 191.
\bibitem{markevitch2007} Markevitch, M. \& Vikhlinin, A. 2007, ``Physics of the intracluster medium'', Phys.\ Rep.\ 441, 1 (cold fronts).
\bibitem{hillier2012} Hillier, A., Isobe, H., Shibata, K. \& Berger, T. 2012, ``Numerical simulations of the magnetic Rayleigh--Taylor instability in the Kippenhahn--Schl\"uter prominence model'', ApJ 746, 120.
\bibitem{hillier2016} Hillier, A. et al. 2016, ``On the nature of the magnetic Rayleigh--Taylor instability in prominences'', MNRAS 462, 2256.
\bibitem{carlyle2017} Carlyle, J. et al. 2017, ``Non-linear growth of the magnetic Rayleigh--Taylor instability in prominences'', A\&A (prep./arXiv).
\bibitem{kalluri2025} Kalluri, M.~T. \& Hillier, A. 2025, ``Self-similarity and growth of non-linear magnetic Rayleigh--Taylor instability --- Role of the magnetic field strength'', Physica D (arXiv:2410.05463).
\bibitem{sefkow2014} Sefkow, A.~B. et al. 2014, ``Simulations of magnetized liner inertial fusion targets...'', Phys.\ Plasmas 21, 072711 (MRT en laboratorio).
\bibitem{boyd2003} Boyd, T.~J.~M. \& Sanderson, J.~J. 2003, \emph{The Physics of Plasmas} (Cambridge Univ.\ Press), p.~99 (campo crítico citado por el \code{init.c}).
\bibitem{ryujones1995} Ryu, D. \& Jones, T.~W. 1995, ``Numerical magnetohydrodynamics in astrophysics: algorithm and tests for one-dimensional flow'', ApJ 442, 228 (tubo de choque MHD de validación).

\end{thebibliography}

\clearpage
\appendix

%=====================================================================
\section{Apéndice A — Números del caso base (cheat-sheet)}
\label{app:numeros}

Configuración \#07 con \code{ETA=2.0}, \code{GRAV=-0.1}, \code{CHI=0.2}, $\gamma=5/3$:

\begin{table}[htbp]
\centering\small
\caption{Todo lo que hay que saber del caso base, ya evaluado. Unidades de código
($L_x=1$, $\rho_1=1$, $c_{s,1}=1$).}
\begin{tabular}{@{}llp{6.9cm}@{}}
\toprule
Magnitud & Valor & Fórmula/origen \\
\midrule
Dominio / malla & $[-\tfrac12,\tfrac12]\times[-1,1]$, $128\times256$ & pluto.ini \#07 \\
$A$ (Atwood) & $1/3$ & $(\eta-1)/(\eta+1)$ \\
$p_0$; $p(y)$ & $0.6$; $0.6-0.1\rho y/1$ & hidrostática; $c_{s,1}=1$ por diseño \\
$c_s$ arriba (en $y=0^+$) & $1/\sqrt2\simeq0.707$ & $\sqrt{\gamma p/\rho_2}$ \\
$B_x$ & $1.784\times10^{-2}$ & $\chi\sqrt{0.1/4\pi}$ \\
$v_A$ abajo / arriba & $1.78\times10^{-2}$ / $1.26\times10^{-2}$ & $B/\sqrt\rho$ \\
$\beta$ (interfaz, ambos lados) & $\simeq3.8\times10^{3}$ & $2p/B^2$ \\
$\sigma_\mathrm{HD}$ & $0.4576$ & $\sqrt{Agk}$, $k=2\pi$ \\
$\sigma$ ($\chi=0.2$) & $0.4484$ ($-2.0\%$) & $\sigma_\mathrm{HD}\sqrt{1-\chi^2}$ \\
$t_\mathrm{e-fold}$ & $2.23$ & $1/\sigma$ \\
$\lambda_c$; $\lambda_\mathrm{max}$ & $0.04$; $0.08$ & $\chi^2L_x$; $2\lambda_c$ \\
$\sigma_\mathrm{max}$ & $1.14$ & $\sqrt{Ag\pi/(2\chi^2L_x)}$ (no sembrado) \\
$\xi_{0,\mathrm{eff}}$ & $\simeq0.022$ & $v_0/(2\sigma)$, $v_0=0.02$ \\
$t$ con $\xi=0.1\lambda$ & $\simeq3.4$ & eq.~\eqref{eq:xi} \\
$t$ con $\xi=\lambda/2$ & $\simeq6.9$ & eq.~\eqref{eq:xi} \\
$v_b$ (Layzer/Goncharov) & $0.126$ & eq.~\eqref{eq:layzer} \\
$\Delta t$ (CFL 0.8) & $\simeq6\times10^{-3}$ & CFL$\cdot\Delta x/c_f$; $\sim$2.5k pasos a $t=15$ \\
Energías iniciales & --- & $E_\mathrm{int}=1.725$;\ \ $E_\mathrm{mag}=3.2\times10^{-4}$;\ \ $E_\mathrm{pot}=-0.05$ (integrales exactas del dominio) \\
\bottomrule
\end{tabular}
\end{table}

%=====================================================================
\section{Apéndice B — Checklist de entregables}
\label{app:checklist}

\begin{itemize}[nosep]
\item[$\square$] \textbf{Informe:} marco teórico (§§\ref{sec:fisica}--\ref{sec:nolineal} con
las ecs.\ 1--13), estado del arte (\S\ref{sec:contexto}+\S\ref{sec:biblio}), descripción del
problema y CI (\S\ref{sec:setup}), metodología (\S\ref{sec:pluto}+\S\ref{sec:python}),
figuras de todas las variables con colorbar/ejes/títulos, análisis comparativo
(\S\ref{sec:comparativo}), conclusiones. Citas completas.
\item[$\square$] \textbf{Código Python:} bien comentado, pipeline completo
(setup $\to$ solver $\to$ postproceso), tests de la escalera de validación incluidos.
\item[$\square$] \textbf{Archivos PLUTO:} \code{definitions.h}, \code{pluto.ini},
\code{init.c} usados + scripts de postproceso (\code{.ipynb} ejecutado de punta a punta,
como en la tarea de Alfvén).
\item[$\square$] Semilla y tiempos idénticos entre códigos; normas $L_1/L_2$ tabuladas;
convergencia discutida.
\end{itemize}

%=====================================================================
\section{Apéndice C — Esqueleto del solver Python}
\label{app:esqueleto}

{\footnotesize
\begin{verbatim}
class RT2D:
    # ---- estado ----
    # U[k] en celdas: rho, mx, my, E       (4, Ny, Nx) float64
    # Bx en caras x: (Ny, Nx+1);  By en caras y: (Ny+1, Nx)
    def __init__(self, Nx, Ny, eta=2.0, g=-0.1, chi=0.2, gamma=5/3):
        ...  # malla, CI eq.(5): p = 1/gamma + rho*g*y (¡discreta! §8.4),
        ...  # semilla (sec. 4.3), Bx = chi*sqrt((eta-1)*|g|)/sqrt(4*pi)
    # ---- conversiones ----
    def con2prim(self): ...        # + energia -> p (con floor numerico)
    def prim2con(self): ...
    # ---- fisica ----
    def cfast(self, prim, Bn): ... # c_f (eq. cfast)
    def hll_flux(self, UL, UR, ...): ...  # flujo HLL
    def reconstruct(self, prim): ...      # MC limiter -> estados L/R
    def emf_corners(self, ...): ...       # CT: promedio de 4 caras
    def source_gravity(self, U): ...      # +rho*g en my ; +rho*(v.g) en E
    # ---- paso ----
    def rhs(self, U, Bx, By): ...  # flujos + CT + fuentes
    def step_rk2(self): ...        # heun: dt por CFL (§8.2)
    def boundaries(self): ...      # x periodico; y reflectivo (receta sec. 3.4)
    # ---- diagnostico ----
    def energy_partition(self): ...    # E_int, E_cin, E_mag, E_pot
    def interface(self): ...           # contorno rho = 1.5 -> y_int(x)
    def divB(self): ...                # monitor ~1e-15 con CT
    def run(self, tstop, dt_out): ...  # loop + guardado npz

# validacion: sigma_hat = polyfit(t, log(a), 1)[0] en 3<=t<=6
# esperado: 0.4484 (chi=0.2) / 0.4576 (chi=0)
\end{verbatim}
}

\end{document}
